\section{Conclusion}

This thesis examines whether the disappearing index effect documented by \citet{greenwood2025} also appears around OMXSB reconstitutions. The Swedish return pattern shows substantial historical attenuation, but no monotonic decline. Additions earn 8.52 \% in 2007--2011 and 1.50 \% in 2022--2026. Across the later decade, the mean is 2.15 \%, with a 95 \% interval of 1.01 to 3.39 \%. The latest block alone is less precise and sensitive to influential events. These findings support a positive later-decade addition association; they establish neither economic disappearance nor a stable current premium.

Reported fund trading and pre-implementation volume confirm that index changes coincide with concentrated benchmark-related activity. In the common later addition sample, covered purchases average 1.03 \% of outstanding shares. The return-to-purchases ratio is 2.63, but changing coverage, broad reporting intervals and potentially shared information prevent a structural interpretation. Regulatory holdings provide a useful check on the later demand scale.

Two extensions sharpen the interpretation. The earliest-to-latest return decline is concentrated after a common initial seven-return window, while the announcement-to-implementation interval becomes shorter. This is a measurement decomposition, not evidence that calendar changes caused the decline. In the selected share-class sample, both classes move substantially and their relative return is imprecise. That pattern is compatible with shared information or demand transmitted across classes; it does not isolate either mechanism.

The Swedish evidence consequently answers the persistence question more clearly than the mechanism question. Positive average addition returns remain in the later sample, alongside observable fund buying, but the data do not identify the causal contribution of that buying. More complete investor coverage, reconciled holdings histories and transaction-level evidence would help distinguish changes in demand from changes in the market's ability to accommodate it.
